\chapter{Introduction}
\label{ch:intro}

\section{The problem}
\label{sec:problem}

A web agent cannot act on a page it has not been shown. Before every click it
must be handed some encoding of the current screen, and there is more than one
way to produce that encoding. The page can be flattened into structured text
drawn from the browser's accessibility tree, the same description of elements,
labels and roles that a screen reader uses. It can instead be rendered as a
screenshot with the candidate elements numbered and boxed on it. Or it can be
sent as both. These are
not equivalent inputs. They differ in what they make visible, and they differ in
what they cost.

The choice is normally made once, at configuration time, and applied to every
step of every task. That is the problem this dissertation is about. A
multi-step web task takes tens of browser steps, and the representation
decision is re-paid at each one. The screenshot is re-encoded, the marks are
re-drawn, and the input is billed again by the token --- the unit in which
models charge for what they are shown. If the extra information changed the outcome at every
step, this would simply be the price of the capability. But it is not obvious
that it does, and, as far as we can establish, nobody has measured whether it does.

The obvious response is to decide per step instead of once. What this
dissertation tests is the weaker, task-level version of that idea: the choice
is made once per task rather than once per step. The reason is that the
benchmark scores whole tasks, so what a different mode would have achieved at
an individual step is never observed --- and that per-step comparison is
exactly what a step-level selector would have to learn from
(Section~\ref{sec:regime-boundary}). That intuition is
what makes this an interesting problem rather than an engineering preference,
because it is only correct if three separate conditions hold together. The cheap and expensive encodings must succeed on \emph{different} tasks, or there is nothing to select between. The need for the expensive one
must be \emph{predictable} from something observable before it is paid for, or
the selection cannot be made in time. And the value recovered must exceed what
the selection machinery itself costs, or the whole apparatus loses money. Each
condition is measurable. This dissertation measures all three, in that order,
and reports where the chain breaks.

\begin{figure}[t]
  \centering
  \includegraphics[width=\textwidth]{fig_f0_thesis_overview.pdf}
  \caption[Thesis overview: the three questions and their answers]{%
  \textbf{The argument of this dissertation in one figure.} The same web task
  can be handed to the same agent through six different page encodings (centre),
  which differ in cost by construction. The questions that follow are answered in that order: \emph{is there a ceiling worth chasing} (yes, but see
  Chapter~\ref{ch:repr} for how much of it a rerun already buys), \emph{is the
  per-task choice predictable at serving time} (no), and \emph{does acting on
  the prediction pay} (the trade-off does not exist under this cost boundary).
  This is a map of the argument, not a performance result; every number in it is
  established in the chapter named beneath it.}
  \label{fig:f0}
\end{figure}

\section{Why the cost is worth arguing about}
\label{sec:stakes}

What turns a per-step configuration choice into a system-level quantity is a combination of three properties of the setting.

\paragraph{The cost is multiplicative in the episode length.} Our agent
averages 13.7--28.4 steps per task depending on site, model and mode
(Chapter~\ref{ch:system}, Table~\ref{tab:cost-per-mode}). A decision that costs
a fraction of a cent per step is a decision made hundreds of thousands of times
across a benchmark sweep, and millions of times in a large deployment.

\paragraph{The expensive part is concentrated, not diffuse.} The annotated-screenshot condition looks expensive because it appears to send more of everything, and that reading is wrong. The
full Set-of-Marks mode is the one that sends a screenshot with the
interactable elements numbered and boxed on it. Measured on matched pages, its text payload
is a median $1.008\times$ that of the structured-text mode on classifieds and $1.004\times$ on reddit (Section~\ref{sec:obs-space};
$n=429$ paired tasks). The cost difference between them is therefore almost
entirely the image. That is a useful fact rather than a curiosity: it means
the expensive component can be removed \emph{surgically}, holding the text and
the prompt fixed, which is exactly the manipulation Chapter~\ref{ch:system}
builds.

\paragraph{Computational cost is not the same thing as environmental cost, and
this dissertation does not conflate them.} We report token counts, billed dollars, elapsed time, and --- for the two
models we run on our own hardware --- the energy the graphics card itself
reports drawing. We do not convert tokens into
kilowatt-hours; no hardware-independent conversion exists
\citep{fernandez2025energy}, and treating one as a proxy for the other is a
known failure mode in this literature \citep{lannelongue2021green,
henderson2020systematic}. Section~\ref{sec:sustainability} states exactly what
our measurement boundary includes and what it therefore cannot support.

\section{Concepts before machinery}
\label{sec:concepts}

The whole argument rests on three ideas, and none of them needs an implementation detail to state.

\paragraph{A page representation is a choice, not a given.} An \emph{observation
mode} is a complete, executable specification of what the agent is shown at a
step: which textual encoding of the page, framed by which prompt, accompanied by
which image, if any. Two modes can expose the same page and still lead the same
model to different actions. Once the representation is a choice rather than a given, each option has a
measurable success rate and a measurable cost, and a choice between options
with measurable outcomes is one that can be made well or badly.

\paragraph{Selection has an upper bound that nobody can reach.} If we look at
outcomes \emph{after} the fact and pick, per task, whichever mode happened to
succeed, we obtain a number that no deployable system can achieve. It is still
worth computing, because it bounds the entire enterprise: it is the most that
perfect selection could ever be worth. We call this the \emph{oracle ceiling},
and we are careful throughout to treat it as a bound on opportunity rather than
as a result.

\paragraph{An upper bound is not a plan.} Between the ceiling and a deployable
system sit two gaps that are usually collapsed in discussion and are kept apart
here. The first is \emph{predictability}: can the right mode be inferred from
signals available before the expensive mode is paid for? The second is
\emph{economics}: even a good predictor loses if the expensive mode is called for too rarely to
matter, if its two kinds of mistake cost different amounts, or if the act of
choosing costs more than the choice saves. The core methodological
commitment of this dissertation is that these three quantities (ceiling, predictability, realised value) are measured separately and never allowed to
stand in for one another.

\begin{figure}[t]
  \centering
  \includegraphics[width=\textwidth]{fig_f1_motivating_example.pdf}
  \caption[What each observation mode actually sends]{%
  \textbf{One page, three ways of paying for it.} All three columns show the same browser step of the same task, verified byte-identical at the source screenshot, as it is actually transmitted to the model. Left:
  the structured page text of the \mDOM{} mode, listing the elements with
  their labels and roles, and no image at all. Centre: the \mSoM{} mode, which
  sends the screenshot with the interactable elements numbered and boxed on it
  \emph{and} the numbered lines of that same text (\sommarks{}). Right: the
  \mVis{} mode, the plain screenshot with no page text. The quantitative point matters more than the pictures: on this
  step the \mSoM{} text payload is $1.008\times$ the \mDOM{} text payload
  (median over all 224 classifieds tasks at step~0: also $1.008\times$), while
  the billed input tokens are $4{,}335$ against $3{,}314$. The extra 31\% is the
  image.}
  \label{fig:f1}
\end{figure}

\section{Approach}
\label{sec:approach}

We construct a controlled representation space over a fixed agent and measure
along the chain just described.

Six observation modes are built by combining three two-way choices: which page
text is supplied (the full structured listing, or only its numbered lines with
the structure flattened away), how the surrounding instructions are worded (as
though the text were the whole input, or as though a numbered screenshot were
coming), and whether the annotated image is sent. Four of the six send no image at all, one sends text and image
together, and one sends only the image. Everything else is held identical
across modes: the task, the model, the settings that govern how the model
produces text, the set of browser actions it may take, the limit on how many
steps it gets, the benchmark's automatic scoring, and the rules for what
counts as a cost. Chapter~\ref{ch:system} makes that claim auditable.

We run this on 8 combinations of website and model, which we call
\emph{cells}. Three models --- two of one design family, one of another ---
are run on two VisualWebArena sites \citep{koh2024visualwebarena}. Two of
those models are then run on a WebArena site \citep{zhou2024webarena}, to
check whether the findings survive in a second benchmark. Every cell runs all six
modes over the same fixed set of tasks, so every comparison is between two
modes attempting the same task.

On that base we then compute four things, in order. First, the ceiling: what
an idealised chooser could reach if it were told after the fact which modes
solved each task, and how far that ceiling falls when each mode in turn is
removed from its menu. Second, a noise floor: how much a success rate moves
when an \emph{unchanged} setup is simply run again, which is the yardstick
every later difference is held against. Third, whether the choice can be
predicted in advance from information a deployed system would have, tested so
that nothing from the tasks being scored leaks into the training. Fourth,
whether acting on those predictions beats fixing one mode in advance, once
cost is counted.

\section{Contributions}
\label{sec:contributions}

\begin{enumerate}[leftmargin=1.5em]
\item \textbf{A controlled decomposition of the page-input choice, and the
measurement that one part of it dominates the cost.} Six working modes
separate the page text, the wording of the instructions, and the presence of
the image, over an otherwise unchanged agent. The factorisation shows that the cost difference between the
      cheap and expensive conditions is almost entirely attributable to the
      image, which is what makes the image-off region worth studying at all
      (Chapter~\ref{ch:system}).

\item \textbf{A ceiling measurement reported next to its own repetition
      baseline.} An idealised after-the-fact chooser gains $+3.45$ to $+16.35$
      percentage points of task success over the best single mode at
      1.6--35.3\% lower cost, in all 8 cells and on both benchmarks. (We write
      \emph{pp} for percentage points from here on.) We pair this with a floor
      measured by rerunning identical setups, and then make the comparison
      fair on count. Adding one plain repeat of an unchanged mode buys
      2.0--7.6\pp{} by itself. Six repeats of the strongest single mode recover
      more than half of everything six different modes add over it, whatever
      the per-task flip rate --- and, if that rate is one half, come within
      3.53\pp{} of matching them outright, a gap smaller than the threshold we
      ourselves derive for what one repeat can produce. Applying that
      standard to our own headline moves the surviving claim from the ceiling's
      height to its price. To our knowledge, no comparable web-agent
      representation study reports its effect sizes against a measured rerun
      floor, let alone against repetition given the same number of attempts
      (Chapter~\ref{ch:repr}).

\item \textbf{A negative prediction result, with the two obvious excuses ruled
out.} We split the tasks repeatedly into a training part and a held-back part.
Which mode to prefer, where to set the decision threshold, and how to scale
the inputs are all decided using the training part alone; only then is the
result measured on the part held back. Under that protocol the learned router
never manages to be \emph{both} more successful \emph{and} cheaper than the
policy that simply always uses the cheapest mode: it clears that joint bar in
\textbf{0 of 8} cells. It does raise success in most cells; what it never does
is raise success without also raising cost. Two controls rule out the
two easiest objections. Shuffling the outcomes against the inputs 10,000 times
shows how large an apparent gain arises from chance alone; and measuring
against always-cheapest, the policy a deployment gets for free, shows that the
shortfall is not an artefact of a weak comparison (Chapters~\ref{ch:pred}--\ref{ch:e2e}).

\item \textbf{A mechanism for the failure, located in label supply rather than
      model capacity.} A task can only say which mode to prefer if some mode
      solved it. At 2--27\% base success, 4 of 6 cells cannot supply two
      classes that each clear the minimum number of training examples we
      require before fitting a classifier at all. Three separate attempts to obtain more labelled
      examples each fail for a different structural reason. We additionally
      show that the strongest single predictor on VisualWebArena is the
      benchmark's own human difficulty annotation, which no deployment possesses. We
      also test the competing ``there simply were not enough examples''
      explanation directly rather than arguing about it. It explains why the
      predictor is weak, but not why the result is negative: the best any
      predictor could become is the idealised chooser, and that chooser mostly
      fails too
      (Chapter~\ref{ch:pred}, Chapter~\ref{ch:disc}).

\item \textbf{A calibrated statement of what the negative result is a negative
      result \emph{about}.} We show that even the idealised chooser fails to
      be both more successful and cheaper than always-cheapest in 7 of 8
      cells. Always-cheapest already has the lowest cost of any fixed policy,
      so any chooser that raises success must also spend more, and the two
      cannot be traded off without an exchange rate this dissertation
      deliberately does not assume. The correct conclusion is therefore that
      no operating point exists that is better on both counts under this cost
      boundary, not that the learner was inadequate (Chapter~\ref{ch:e2e}).
\end{enumerate}

\section{This dissertation contains a ``no'', and that is deliberate}
\label{sec:the-no}

The second half of the research question (can the need be predicted cheaply enough to be worth predicting?) is answered in the negative. Two consequences
follow, and both are structural rather than cosmetic.

First, Chapter~\ref{ch:e2e} is not organised around how much the router won. It
is organised around \emph{where the chain breaks and why that break is a
property of the regime rather than of our implementation}. Second, the
scientific contribution lives in Chapter~\ref{ch:disc}: what this work delivers
is not a router but a boundary: a characterisation of the conditions under
which representation routing cannot be learned.

A negative result can carry a dissertation only if two things are established
first, and we take both as obligations rather than as options. The opportunity
must be shown to exist, or ``we could not learn it'' collapses into ``there was
nothing there''; this is what the ceiling measurement in Chapter~\ref{ch:repr}
is for. And the failure must be traced to a named mechanism, or ``we could not
learn it'' collapses into ``your method was bad''; this is what the label-supply
diagnosis in Chapter~\ref{ch:pred} is for. Where our evidence stops short of
either obligation, we say so in place rather than in a closing list of caveats.

\section{What this dissertation does not claim}
\label{sec:not-claimed}

The statements below are consistent with our data but stronger than it, so we do not make them anywhere in this document.

\begin{itemize}[leftmargin=1.5em]
\item That the image-off modes \emph{replace} the full Set-of-Marks mode. Their
      interest is complementarity, and Chapter~\ref{ch:repr} shows even that is
      weaker than an earlier reading of our own data suggested.
\item That dropping the image costs little performance, or that the image-off
      configuration constitutes a cheaper Set-of-Marks. On four of six
      VisualWebArena cells the full mode is the strongest single mode.
\item That representation routing is unlearnable in general. We characterise why
      it is unlearnable \emph{in the observed 2--27\% base-success regime}, and
      Section~\ref{sec:regime-boundary} states what would have to change for the
      conclusion to change.
\item That reduced token or dollar cost constitutes a reduced environmental
      footprint. Section~\ref{sec:sustainability} names the exact measurement
      boundary and what it excludes.
\end{itemize}

\section{Map of the dissertation}
\label{sec:map}

Each chapter answers one question and hands the next chapter the question it
could not answer.

\begin{description}[leftmargin=0em,itemsep=3pt]
\item[Chapter~\ref{ch:bg}] asks what prior work already settles, and shows that
  four active research programmes each optimise a different decision variable
  than the one studied here.
\item[Chapter~\ref{ch:system}] asks how representation choice can be compared
  fairly enough for later differences to be interpretable, and fixes the
  comparison boundary that the rest of the document depends on.
\item[Chapter~\ref{ch:repr}] asks whether the modes differ enough for selection
  to be worth anything, establishes the ceiling, and then measures how much of
  it survives contact with a rerun.
\item[Chapter~\ref{ch:pred}] asks whether the need can be predicted from
  serving-time signals, and separates ``predictable'' from ``useful''.
\item[Chapter~\ref{ch:e2e}] asks whether acting on those predictions beats a
  fixed policy once cost is counted, and finds that the comparator itself
  explains most of the answer.
\item[Chapter~\ref{ch:disc}] reconciles the three results, states the regime in
  which they hold, and closes on what the work enables.
\end{description}

\chapterhandoff{This chapter has argued that choosing how to show the page ---
rather than fixing that choice once --- is an efficiency problem with
measurable structure, and has set out the three conditions under which it
would be worth solving. It has not shown that the question is still open.
Chapter~\ref{ch:bg} therefore examines four neighbouring lines of work: what a
web agent is shown, choosing between models, letting a model decline to answer
when unsure, and varying how much computation a case receives. Each optimises
a different decision than the one studied here.}
